\section{Discussion and Conclusion}

In this paper, we have shown how PTD-P (inter-node pipeline parallelism, intra-node tensor parallelism, and data parallelism) can be composed to achieve high aggregate throughput (502 petaFLOP/s) while training large models with a trillion parameters. This facilitates end-to-end training in reasonable times (estimated time of around 3 months for a trillion-parameter model). We discussed the various tradeoffs associated with each of these types of parallelism, and how the interactions between them need to be considered carefully when combined.

Even though the implementation and evaluation in this paper is GPU-centric, many of these ideas translate to other types of accelerators as well. Concretely, the following are ideas that are accelerator-agnostic: a) the idea of smartly partitioning the model training graph to minimize the amount of communication while still keeping devices active, b) minimizing the number of memory-bound kernels with operator fusion and careful data layout, c) other domain-specific optimizations (e.g., scatter-gather optimization).